% chapters/chapter6_discussion.tex
\chapter{Discussion}
\label{ch:discussion}

% ============================================================
% CHAPTER OVERVIEW
% ============================================================
% This chapter interprets the results:
% - What is the dominant source of discriminability?
% - What does this mean for Apple Catcher deployment?
% - What does this mean for visual MI-BCI more broadly?
%
% TARGET: 10-15 pages total
%
% STYLE NOTES:
% - Now interpret, don't just describe
% - Connect findings to hypotheses
% - Discuss implications
% - Be honest about limitations


% ============================================================
% 6.1 PRIMARY SOURCE OF DISCRIMINABILITY
% ============================================================
\section{What Drives Classification in Apple Catcher?}
\label{sec:discuss_primary}

% TARGET: 4-5 pages


% ------------------------------------------------------------
\subsection{Evaluating the Competing Hypotheses}
\label{subsec:discuss_hypotheses}

% TARGET: 2-2.5 pages
%
% CONTENT:
%
% Go through each hypothesis and synthesize the evidence:
%
% ARTIFACT HYPOTHESIS:
% - Summarize evidence for: [list tests that supported]
% - Summarize evidence against: [list tests that did not support]
% - Overall assessment: [strong/moderate/weak support]
%
% If artifacts play a significant role:
% - This is concerning for deployment
% - But not necessarily catastrophic (see Section 6.3)
% - Discuss what type of artifacts (eye movement? muscle?)
%
% ERD HYPOTHESIS:
% - Summarize evidence for and against
% - Did ERDS show expected patterns?
% - Do filters and attributions align with ERDS?
%
% If ERD contributes significantly:
% - This is the "good" outcome
% - Model learns legitimate MI signal
% - Validates the paradigm
%
% MRCP HYPOTHESIS:
% - Summarize evidence
% - Was MRCP present in the data?
% - Do low-frequency filters and early attribution support MRCP?
%
% If MRCP contributes:
% - Interesting finding about motor preparation
% - Explains pre-cue superiority
% - Has implications for paradigm design
%
% VISUAL HYPOTHESIS:
% - Summarize evidence
% - Was occipital activity present and discriminative?
%
% If visual processing contributes:
% - Model may classify based on visual attention, not MI
% - Less concerning than artifacts but still problematic


% ------------------------------------------------------------
\subsection{The Most Likely Explanation}
\label{subsec:discuss_explanation}

% TARGET: 1.5-2 pages
%
% CONTENT:
%
% Synthesize into a coherent explanation:
%
% SCENARIO A: Artifacts Dominate
% If evidence points primarily to artifacts:
% - The 78% accuracy is largely or partially invalid
% - Pre-cue superiority explained by eye tracking
% - Model would fail for paralyzed patients
% - Recommendation: artifact rejection or paradigm redesign
%
% SCENARIO B: ERD Dominates
% If evidence points primarily to ERD:
% - The 78% accuracy reflects legitimate MI classification
% - Pre-cue ERD may reflect motor preparation
% - Model should generalize to patients
% - Recommendation: proceed with deployment
%
% SCENARIO C: MRCP Dominates
% If evidence points primarily to MRCP:
% - Interesting finding about anticipatory processing
% - Model learns motor preparation, not sustained imagery
% - May or may not generalize to patients
% - Recommendation: investigate further
%
% SCENARIO D: Mixed Sources
% Most likely outcome - multiple contributions:
% - Estimate relative contributions if possible
% - E.g., "60% ERD, 25% MRCP, 15% artifact"
% - Discuss which dominates pre-cue vs post-cue
%
% BE HONEST:
% - If you cannot definitively distinguish, say so
% - Note what additional data/analysis would be needed


% ------------------------------------------------------------
\subsection{Explaining the Pre-Cue Superiority}
\label{subsec:discuss_precue}

% TARGET: 1 page
%
% CONTENT:
%
% The central finding from prior work:
% - Pre-cue window (-1 to 0s) alone: 70.7%
% - Post-cue window (0 to 3s) alone: 58.0%
% - This is counterintuitive for standard MI-BCI
%
% Possible explanations given your findings:
%
% If artifacts:
% - Eye movements tracking falling apple
% - Concentrated in pre-cue period
% - Explains the pattern well
%
% If MRCP:
% - Motor preparation before cue
% - Bereitschaftspotential or CNV
% - Makes sense given paradigm design
%
% If ERD with early onset:
% - Subjects start imagery when they see apple
% - Don't wait for formal cue
% - ERD begins early, explaining pre-cue power
%
% The interpretation depends heavily on the results you found.


% ============================================================
% 6.2 IMPLICATIONS FOR APPLE CATCHER DEPLOYMENT
% ============================================================
\section{Implications for Apple Catcher System}
\label{sec:discuss_implications_ac}

% TARGET: 3-4 pages


% ------------------------------------------------------------
\subsection{If Artifacts Contribute Significantly}
\label{subsec:discuss_if_artifacts}

% TARGET: 1-1.5 pages
%
% CONTENT:
%
% THE CONCERN:
% - Stroke patients (target users) may not produce same artifacts
% - Paralysis affects not just motor but also oculomotor control
% - Model trained on artifacts won't help these patients
%
% SEVERITY ASSESSMENT:
% - If artifacts are primary: system needs fundamental redesign
% - If artifacts are secondary: mitigation may be sufficient
%
% MITIGATION OPTIONS:
%
% Option 1: Artifact Rejection
% - Add EOG channels for monitoring
% - Reject trials with eye movements
% - Cost: lose some trials, adds hardware
%
% Option 2: Artifact-Aware Training
% - Include artifact-contaminated and clean trials
% - Model learns to rely on neural signal
% - Requires labeled artifact data
%
% Option 3: Paradigm Redesign
% - Make visual cue less attention-grabbing
% - Central cue instead of lateralized apple
% - Loses gamification benefit
%
% Option 4: Post-Cue Only
% - Use only post-cue window
% - Accept lower accuracy (58%)
% - Avoids artifact-heavy period


% ------------------------------------------------------------
\subsection{If Neural Signal Dominates}
\label{subsec:discuss_if_neural}

% TARGET: 1-1.5 pages
%
% CONTENT:
%
% VALIDATION:
% - The 78% accuracy is scientifically meaningful
% - Model learns legitimate motor imagery signal
% - System can proceed to clinical validation
%
% NEXT STEPS:
% - Test on actual stroke patients
% - Validate generalization
% - Compare to traditional BCI systems
%
% THE PRE-CUE FINDING AS OPPORTUNITY:
% - If MRCP or early ERD explains pre-cue power
% - This is actually useful information
% - Motor preparation signals may be more robust
% - Could inform future paradigm design


% ------------------------------------------------------------
\subsection{Recommendations for Deployment}
\label{subsec:discuss_recommendations}

% TARGET: 1 page
%
% CONTENT:
%
% Concrete recommendations based on findings:
%
% If primarily artifacts:
% 1. Do not deploy without mitigation
% 2. Add EOG monitoring
% 3. Consider paradigm redesign
% 4. Revalidate with artifact-free data
%
% If primarily neural:
% 1. Proceed to patient testing
% 2. Keep pre-cue window (it works)
% 3. Monitor for artifact development in patients
%
% If mixed:
% 1. Quantify contributions more precisely
% 2. Implement artifact monitoring
% 3. Cautious deployment with tracking


% ============================================================
% 6.3 BROADER IMPLICATIONS FOR VISUAL MI-BCI
% ============================================================
\section{Implications for Visual Motor Imagery Paradigms}
\label{sec:discuss_implications_broad}

% TARGET: 2-3 pages


% ------------------------------------------------------------
\subsection{The Visual Cue Confound Problem}
\label{subsec:discuss_visual_confound}

% TARGET: 1-1.5 pages
%
% CONTENT:
%
% GENERALIZABLE LESSON:
% - Apple Catcher is not unique
% - Any MI-BCI with visual lateralized cues has this problem
% - Examples: left/right arrows, spatial targets
%
% The issue:
% - Visual stimulus can create:
%   * Eye movements (artifact)
%   * Visual attention shifts (neural but not MI)
%   * Visual evoked potentials
% - All of these can be lateralized and class-discriminative
%
% RECOMMENDATION FOR FIELD:
% - Always check for these confounds
% - Interpretability analysis should be standard
% - Don't assume high accuracy = legitimate signal


% ------------------------------------------------------------
\subsection{The Need for Interpretability in BCI}
\label{subsec:discuss_need_interpretability}

% TARGET: 1-1.5 pages
%
% CONTENT:
%
% THE PROBLEM:
% - BCI literature often reports accuracy without validation
% - Deep learning makes this worse (black box)
% - High accuracy can mask artifact exploitation
%
% THE SOLUTION:
% - Interpretability should be standard practice
% - Compare learned features to expected neurophysiology
% - Check for artifact signatures
%
% THIS THESIS AS EXAMPLE:
% - Demonstrates the approach
% - Shows what questions to ask
% - Provides methodology template


% ============================================================
% 6.4 LIMITATIONS
% ============================================================
\section{Limitations}
\label{sec:discuss_limitations}

% TARGET: 2-3 pages


% ------------------------------------------------------------
\subsection{Dataset Limitations}
\label{subsec:discuss_limit_dataset}

% TARGET: 0.5-1 page
%
% CONTENT:
%
% Single dataset:
% - All findings specific to Apple Catcher
% - May not generalize to other MI paradigms
%
% Healthy subjects only:
% - No stroke patients in dataset
% - Cannot validate clinical relevance
%
% Unknown apple timing:
% - Exact onset of apple falling is undocumented
% - Limits precision of pre-cue analysis
%
% No dedicated EOG channels:
% - Cannot definitively separate eye artifacts
% - Frontal channels used as proxy


% ------------------------------------------------------------
\subsection{Methodological Limitations}
\label{subsec:discuss_limit_methods}

% TARGET: 0.5-1 page
%
% CONTENT:
%
% Offline analysis:
% - All analysis done post-hoc
% - Online behavior may differ
% - No closed-loop validation
%
% Single model:
% - Analyzed only baseline EEGNet
% - Fine-tuned/DANN models may differ
% - Could be explored in future work
%
% Attribution method choice:
% - Used Integrated Gradients
% - Other methods might give different results
% - IG has known limitations


% ------------------------------------------------------------
\subsection{Interpretation Limitations}
\label{subsec:discuss_limit_interpretation}

% TARGET: 0.5-1 page
%
% CONTENT:
%
% Correlation vs causation:
% - Attribution shows what's used, not necessarily why
% - Cannot prove causal mechanism
%
% Cannot fully resolve hypotheses:
% - Some ambiguity likely remains
% - Multiple sources may contribute
% - More targeted experiments needed for certainty
%
% Alternative explanations:
% - May have missed relevant hypotheses
% - Unknown confounds could exist


% ============================================================
% 6.5 FUTURE DIRECTIONS
% ============================================================
\section{Future Directions}
\label{sec:discuss_future}

% TARGET: 1-2 pages
%
% CONTENT:
%
% IMMEDIATE FOLLOW-UPS:
%
% 1. EOG-controlled validation
%    - Collect new data with EOG channels
%    - Definitively separate eye artifacts
%    - Re-run interpretability analysis
%
% 2. Patient validation
%    - Test model on stroke patients
%    - Compare performance to healthy subjects
%    - Check if same features are used
%
% 3. Paradigm modifications
%    - Test with non-lateralized cues
%    - Compare interpretability results
%    - Isolate visual confound contribution
%
% METHODOLOGICAL EXTENSIONS:
%
% 4. Analyze other models
%    - Compare DANN vs fine-tuned vs baseline
%    - Do different models learn different features?
%
% 5. Causal validation
%    - Ablation studies: remove channels, measure impact
%    - Frequency band filtering experiments
%
% 6. Alternative attribution methods
%    - Compare IG to SHAP, LRP, GradCAM
%    - Check for consistency


% ============================================================
% 6.6 CHAPTER SUMMARY
% ============================================================
\section{Chapter Summary}
\label{sec:discuss_summary}

% TARGET: 0.5 page
%
% CONTENT:
%
% Key takeaways:
% 1. [Main finding about what drives classification]
% 2. [Implication for Apple Catcher deployment]
% 3. [Broader lesson for visual MI-BCI]
% 4. [Limitations to keep in mind]
